\minitopic{运气的类型学}事件运气指某事发生本身偶然，例如一阵风把文件吹到研究者脚边；能力运气指主体碰巧拥有良好教育或天赋；证据运气指主体偶然遇到关键资料；真值运气则指信念以当前方式形成时很容易为假。反运气知识论主要排除最后一种，而非要求认识生活完全没有幸运。 这种区分避免过强结论。科学发现常受偶遇启发，但研究者随后验证，最终信念不再容易为假。一个人幸运出生在拥有可靠学校的环境，也不因此从未学到知识。评价焦点是，目标信念之真是否以适当方式依赖认识能力和证据。 还有\term{介入运气}{intervening luck}与\term{环境运气}{environmental luck}。前者在信念形成链中插入干扰并碰巧抵消，如磁场射击；后者的实际过程本身正常，但周围充满相似错误机会，如假谷仓。二者对能力归因的影响可能不同：环境运气中主体确实准确运用视觉，却仍不安全。

\minitopic{模态条件的精确读法}敏感性检验以$p$为假的最近世界为起点，询问主体是否仍相信$p$；安全性以主体仍以相关方式相信$p$的最近世界为起点，询问$p$是否仍真。二者条件化方向不同，不能用“在其他世界会不会错”笼统代替。 设主体看到正常工作的仪表显示油箱有油。若油箱空，仪表会显示空，信念敏感；在主体按仪表相信有油的近邻世界，通常确有油，信念安全。若仪表卡住在“有油”，实际恰好有油：在空油箱世界仍显示有油，不敏感；主体以此方式相信的近邻世界常为假，也不安全。 彩票案例更复杂。相信“我的票会输”通常敏感：在票中奖的最近世界，开奖信息若尚未公布，主体仍会相信会输，所以其实不敏感；若信念是在核对开奖结果后形成，则敏感且安全。形成基础必须写入反事实评估，不能只看命题概率。

\minitopic{安全性与闭合}安全性常被认为比敏感性更容易保留闭合，但闭合仍非自动。$p$安全、$p$蕴涵$q$，并不在所有描述下保证$q$信念使用同一安全基础。主体可能通过复杂推理引入错误，或$q$本身把遥远怀疑情景变成相关。 更有限的原则是：若主体安全地知道$p$，胜任地认识$p$蕴涵$q$，并以安全的演绎过程推出$q$，则$q$可成为知识。每个限定都有工作：胜任条件排除逻辑误算，基础条件防止主体因另一不可靠理由接受$q$，模态条件确保推理没有把错误机会放大。 怀疑论命题“我不是缸中脑”仍造成压力。若主体在近邻世界中以相同经验相信它，而该世界为缸中脑，信念不安全。安全理论可说缸中脑世界并不近，或者区分普通知觉基础与显式反怀疑推理。接近性判断在这里承担实质理论负担。

\minitopic{安全是否允许风险阈值}绝对安全要求所有相关近邻信念都真，可能过强。现实能力总有极低故障率。概率化安全把要求改为：在相关模态区域内，以该方式形成的信念为真比例足够高。但这又接近可靠主义，并引入阈值问题。 阈值不必是单一数字。证据种类、错误结构和是否存在个案联系可能比总概率重要。正常知觉的统计可靠率未必可精确给出，却有对象持续造成经验的联系；彩票命题概率极高，却只有对称统计，没有把这张票与失败结果区分开的个案信息。 一种解释说，知识要求把主体置于事实的“信息通道”，使信念能随事实变化。概率是通道质量的一项指标，而非完整定义。这样可以解释为什么99.999\%的彩票预测与较低统计准确率的对象识别获得不同知识判断。

\minitopic{德性认识论的两条路线}可靠主义德性认识论把德性理解为稳定产生真信念的能力，如视觉、记忆和良好推理。责任主义德性认识论把开放、认真、勇气、谦逊和公正等人格品质置于中心\parencite{zagzebski1996}。前者更直接分析知识，后者更适合评价探究者与教育目标。 能力成功理论常使用“因为”或归功关系：主体因为能力而成功，而不仅是能力参与了成功。射手的技术参与磁场抵消案例，却不是命中的充分解释。认识成功也可能由能力、环境和他人共同解释；“归功于主体”若要求单独功劳，会排除证言和团队知识。 反运气德性理论尝试结合：信念在近邻情形不易错，并且其真在显著程度上归功于主体的认知能力\parencite{pritchard2012}。结合处理环境运气与能力成就，却可能使两个条件互相重叠或在困难案例中给出冲突判定。

\minitopic{反射知识与认识能动性}一阶能力可以无须主体说明机制。优秀球员接球、儿童辨认父母、动物找到食物，都可能表现可靠认知。反射知识则要求主体对自己如何知道有适当视角：能识别来源、相关风险与能力边界。 反思不能要求无限层级。若知道自己知道还需知道自己知道自己知道，回溯不会终止。合理要求是，主体在当前任务上能响应显著击败者，并以不过度幸运的方式评价一阶能力。医生无需证明整个医学方法，但应知道检测适用人群、误差和替代诊断。 认识能动性也包括主动塑造环境：选择光线更好处观察、保存来源、请独立者复核。知识不是只在头脑内部筛选信念，主体可以通过行动让信念更安全。制度和工具因而能够成为德性表现的一部分。

\minitopic{能力—环境的生态观点}能力没有脱离环境的固定可靠度。利用地标导航在熟悉城市有效，在沙漠中失败；相信长者在稳定小共同体可能获得信息，在大规模操纵网络中易受骗。生态观点研究主体策略与环境信息结构的配合。 环境也可被设计成认识友好或敌对。标准化标签、审计记录与透明更正降低错误机会；深度伪造、暗黑模式和选择性呈现增加环境运气。要求个人以更强意志抵抗所有操纵不现实，认识规范应同时评价平台和制度。 德性教育于是有双重任务：培养辨别、推理和谦逊能力，也教会主体建构支持这些能力的环境。关闭干扰、比较独立来源、使用检查清单，不是能力不足的拐杖，而是高阶认识能力的表现。

\minitopic{综合案例：实验室自动化}实验机器人通常准确移液。某天吸头堵塞，使样本量不足；传感器也恰好读低并触发补液，最终总量正确。研究者根据系统日志相信浓度合格。结果真，却由两项故障抵消。 介入运气明显破坏个案安全。如果实验室有独立称重并在发布前验证，研究者的新信念可基于验证获得知识。若只有系统长期99.9\%准确率，总体可靠不一定修复个案路径。若研究者不知道故障也无从发现，他可能无责，却仍缺少知识。 团队如何改进？记录原始传感器数据、对关键步骤使用异质冗余、把异常抵消列为风险、定期注入故障测试。这样的工程措施把反事实分析转成实践：系统在一处变化时能否响应，错误是否会被独立通道发现。安全性不再只是可能世界术语，也是鲁棒设计原则。

\begin{tcolorbox}[breakable,colback=white,colframe=blue!30!black,title={反运气比较表},fonttitle=\bfseries]
对每个案例分别回答：若$p$为假还会信吗；以同样方式相信时容易假吗；成功是否归功于稳定能力；运气介入链条还是只存在于环境；主体有无反思与改造环境的能力；加入何种独立核验可以产生知识。
\end{tcolorbox}
